\section{Process、Thread 与内核服务总览}

\topic{进程是资源所有权集合}
进程包含地址空间、线程、打开文件、权限和生命周期状态。线程保存可调度执行
上下文。调度器在时钟中断或阻塞点选择可运行线程，并保存/恢复寄存器与栈。
抢占发生在任意允许中断的位置，因此共享状态必须有同步规则。

\topic{锁、等待与 IPC}
自旋锁适合极短且不能睡眠的临界区；等待队列把不能继续的线程从运行队列移出；
管道把字节缓冲、读写端引用和唤醒条件组合成 IPC 对象。任何同步原语都要说明
谁能睡眠、谁能唤醒、锁顺序和中断上下文约束。

配套 v0.10 已实现其中可独立验收的最小闭环：原子 acquire/release 自旋锁、
\code{Blocked} 与具名等待原因、Try/Wait 分离的系统调用，以及 64 字节单
生产者/单消费者管道。生产者写入 256 字节确定性序列，消费者用 31 字节缓冲
逐字节校验；容量和块长不整除，迫使环形回绕、部分传输、满/空阻塞与唤醒实际
发生。通用等待队列容器、可睡眠 mutex 和多端点引用仍属于后续阶段，正文
不会把固定 bootstrap pipe 误称为完整 POSIX 管道。

\topic{从块到文件}
块设备只认识扇区。文件系统用超级块、inode、目录项和数据块提供名字与层次；
VFS 用统一对象接口连接进程文件描述符和具体文件系统。写操作还需要区分：
复制进缓存、提交给设备、设备报告完成、数据达到持久介质。

配套 v0.11 已实现其中刻意有界的完整闭环：2 MiB 磁盘把前 1 MiB 固定给
启动链，后续 1024 个扇区组织为 superblock、bitmap、32 个 inode、目录项和
数据块；八项 LRU 写回缓存最终通过 ATA PIO write 与 FLUSH CACHE 持久化。
每个 PCB 有四个普通文件句柄，Ring 3 可执行 Open/Read/Write/Close/Mkdir/
Sync。当前尚无把管道、文件和设备统一起来的 VFS 打开对象，因此“已有 fd”
不等于“已有完整 VFS”。

\topic{失败与恢复}
磁盘可能在任意写入之间掉电。更新多个结构时，需要日志、写时复制或明确的
恢复算法保证旧状态或新状态至少一个完整。v0.11 的 Dirty/Clean 协议只保证
下次挂载能识别未完成事务并停止，不保证重放恢复；这是显式的阶段边界，而
不是日志文件系统。测试同时覆盖跨块文件、重挂载、随机改写、孤儿 inode、
非法目录名、同盘双启动和超级块损坏，不把创建一个小文件当作完成。

\topic{程序、进程与线程的区别}
程序是磁盘上的代码与数据格式，进程是一次运行实例及其资源所有权，线程是
进程内可被调度的执行流。同一程序可以产生多个进程，一个进程也可以包含多个
线程。把三者分开后，ELF 装载、地址空间、文件描述符和寄存器现场才有清楚的
归属。

进程控制块保存标识、状态、地址空间、线程集合、打开文件和退出信息。线程控制
块保存通用寄存器、内核栈、调度状态和等待原因。对象之间通过稳定标识和受控
引用关联，不能依靠生命周期不明的裸指针。

\topic{上下文切换保存继续执行所需的最小状态}
调度器选中另一个线程后，汇编切换例程保存被调用者保存寄存器和栈指针，把
当前线程状态置入运行队列，再加载目标线程的栈与寄存器。地址空间变化时还要
切换 CR3；用户任务需要恢复 FS/GS 基址及后续扩展状态。

不必在每次切换中保存由调用约定允许破坏的寄存器，因为调度发生在统一入口，
上层中断桩已经保存完整用户现场。保存集合由两层 ABI 共同决定，必须用汇编
偏移断言与结构布局测试保持一致。

\topic{调度从策略中分离机械动作}
运行队列维护可运行线程，阻塞线程位于等待队列，当前线程占用处理器。最简单
的轮转策略按固定时间片选择下一个线程，能够验证抢占和公平推进；优先级、
多核负载均衡和实时期限属于后续策略。

机械切换只负责保存与恢复，调度策略只选择对象。这样可以在不触碰脆弱汇编
上下文的前提下更换算法，也能分别测试状态机和选择顺序。

\topic{睡眠与唤醒必须避免丢失事件}
线程等待条件时，不能先释放锁再随意把自己标为睡眠；事件可能恰好发生在两个
动作之间，导致线程永久错过唤醒。正确协议在持锁状态下检查条件并加入等待队列，
再原子地释放锁与阻塞。唤醒方在同一同步规则下修改条件并移动等待者。

等待通常写成循环而不是一次判断，因为唤醒只表示条件可能变化，资源可能已被
另一个线程取得。中断、超时和进程终止还会形成额外唤醒原因，返回状态必须让
调用者区分。

当前实现利用单核 \code{INT 0x80} interrupt gate 清 IF 的窗口，让再次检查
条件与 PCB 进入 Blocked 不被本地普通 IRQ 插入；管道字段本身由自旋锁保护。
这不是 SMP 完成态：另一 CPU 不受本地 IF 约束，未来必须让条件状态和等待者
登记共同受等待队列锁保护。

\topic{文件描述符把整数映射到内核对象}
用户进程看到的小整数是进程文件表中的槽位，不直接表示磁盘位置。槽位引用打开文件
描述，其中包含对象类型、当前位置、访问模式和引用计数。多个描述符可以通过
复制共享同一打开描述，因而共享文件偏移；分别打开则拥有独立偏移。

VFS 把普通文件、目录、设备和管道组织为统一操作集合，但各对象仍有不同语义。
\code{read} 对文件可能遇到 EOF，对终端可能阻塞，对管道还依赖写端是否存在。
统一接口只统一调用形式，不抹去对象状态。

\topic{inode 与目录把名称连接到数据}
inode 保存文件类型、长度、权限和数据块映射，不保存自身路径。目录是一种从
名称到 inode 编号的映射；路径解析从根或当前目录逐段查找。硬链接让多个目录
项引用同一 inode，引用计数归零且没有打开引用后，数据块才可回收。

磁盘格式中的所有字段使用明确端序和宽度，不能直接把宿主 C++ 结构体写入扇区。
结构体可能有填充、对齐和 ABI 差异。编码器逐字段读写，并在使用长度、块号和
偏移前验证范围。

\topic{缓存把性能问题变成一致性问题}
缓冲缓存保存最近访问的磁盘块，减少慢速 I/O。脏块表示内存内容尚未持久化，
回写线程按策略提交。缓存对象需要引用计数、锁和状态转换，防止一个线程读取
另一个线程正在填充的半成品。

用户看到写调用返回，并不必然表示数据已经到达稳定介质。接口要区分写入页缓存、
提交到设备与设备持久完成；\code{fsync} 一类操作才请求更强保证。QEMU 镜像
仍能用于断电点故障注入，验证恢复算法。

\topic{日志把多块更新变成可恢复事务}
创建文件可能同时修改目录、inode 位图、inode 和数据位图。若写到一半断电，
磁盘会处于跨结构不一致状态。预写日志先把将要发生的更新记录并持久化，再写
正式位置；恢复时根据提交记录重放完整事务或忽略未提交事务。

日志解决崩溃一致性，不自动解决并发隔离、数据校验或硬件撒谎。实现要定义事务
大小、日志环绕、写入顺序和屏障，并用每个可能断电边界的镜像重启测试验证。

\topic{Shell 是整个系统契约的汇合点}
Shell 读取终端输入，解析命令，创建进程，连接管道和文件描述符，等待退出状态。
一个简单命令会同时经过键盘中断、终端缓冲、系统调用、VFS、ELF 装载、地址
空间和调度器。它适合作为 v1.0 的端到端验收，却不能代替各层独立测试。

基础命令保持小而可解释，先覆盖输出、文件查看、目录和进程等待。用户态程序
使用独立 ABI 头文件与最小运行库，不链接宿主 libc，从而确保所有服务确实由
自研内核提供。
